\volpart{第三编\quad 几何相位与现代量子多体（8.513）}
\originalchapter{绝热几何、Chern数与波包}
\label{ch:berry}
\source{8.513第1--6讲幻灯片，结合8.511能带与Hall效应。}
\originalsection{Berry相位的推导}
缓慢参数$R(t)$使$H(R)\ket{n(R)}=E_n(R)\ket{n(R)}$。若能级孤立且$|\hbar\avg{m|\dot H|n}|\ll|E_m-E_n|^2$，绝热态可写$\ket{\psi(t)}=\e^{-(\ii/\hbar)\int E_n\dd t}\e^{\ii\gamma_n(t)}\ket{n(R(t))}$。代回Schrödinger方程后动力学相位抵消，剩$\dot\gamma_n=\ii\avg{n|\partial_t n}$。因此$\gamma_n=\int\mathcal A\cdot\dd R$，$\mathcal A=\ii\avg{n|\nabla_Rn}$。

本征态相位$\ket n\to\e^{\ii\zeta(R)}\ket n$使$\mathcal A\to\mathcal A-\nabla\zeta$。闭合路的$\e^{\ii\gamma}$不变，开路积分则要同时指定端点相位。Berry联络是选局部相位规范后的几何对象，曲率$\mathcal F_{ij}=\partial_i\mathcal A_j-\partial_j\mathcal A_i$规范不变。用本征方程微分，可把它写成
\begin{equation}
\mathcal F_{ij}=-2\operatorname{Im}\sum_{m\ne n}\frac{\avg{n|\partial_iH|m}\avg{m|\partial_jH|n}}{(E_n-E_m)^2},
\end{equation}
此式适合计算，且说明能隙闭合附近曲率为何可能变大。
\originalsection{两能级系统与拓扑整数}
$H=d_0I+\vct d\cdot\vct\sigma$的本征值$E_\pm=d_0\pm|d|$。用球角表示$\hat d$，上能级可选$\ket+= (\cos(\theta/2),\e^{\ii\phi}\sin(\theta/2))^T$，则$\mathcal A_\phi=-\sin^2(\theta/2)$、$\mathcal F_{\theta\phi}=-\sin\theta/2$；下能级符号相反。闭合路径的相位为带符号立体角的一半。

二维绝缘体占据带的曲率积分定义$C=(2\pi)^{-1}\int_{\mathrm{BZ}}\mathcal F_{k_xk_y}\dd^2k$。为理解整数性，把周期Brillouin环面覆盖为局部规范片；每片上由Stokes定理曲率积分化为边界联络积分，相邻片的差为$-\dd\zeta$。相位过渡函数$\e^{\ii\zeta}$沿闭合边界只能有整数绕数，所以总积分是$2\pi$的整数。$C\ne0$时无法在整个环面选一组既光滑又周期的单带本征矢。

对于两带模型，下带$C=(4\pi)^{-1}\int\hat d\cdot(\partial_{k_x}\hat d\times\partial_{k_y}\hat d)\dd^2k$，在本章联络约定下取这个正号。它是从Brillouin区到球面的带向映射次数。连续改变参数且保持$|d|>0$不能改变整数，因此相变必须闭隙或破坏所要求的对称性。单个连续Dirac锥可贡献半整数的局部曲率积分，完整晶格的其他动量区使总数为整数，不能把未经正则化的一只锥当成独立二维晶格绝缘体。
\originalsection{几何修正的运动方程}
孤立带波包的作用量取
\begin{equation}
L=(\hbar\vct k-e\vct A)\cdot\dot{\vct r}+\hbar\mathcal A_n(\vct k)\cdot\dot{\vct k}-E_n(\vct k)+e\phi.
\end{equation}
对$\vct r,\vct k$变分得到$\hbar\dot{\vct k}=-e(\vct E+\dot{\vct r}\times\vct B)$、$\dot{\vct r}=\hbar^{-1}\nabla_kE_n-\dot{\vct k}\times\vct\Omega_n$。反常速度在电场下给横向电流；填满二维带时普通群速度积分为零，而曲率积分给$\sigma_{xy}=Ce^2/h$，整体符号还随Hall电流与场方向定义。此推导把拓扑整数与可测输运相连。

有磁场时相空间体积也修正为$[1+(e/\hbar)\vct B\cdot\vct\Omega]$倍，因运动坐标不再有简单的正则辛结构。对态密度、磁化与输运做积分时要保留这项；只改速度可能违反连续性。强带间混合、闭隙点及极强外场超出单带波包近似。
\originalsection{拓扑泵与边界}
一维绝缘体若参数随时间周期变化，将$(k,t)$视为二维参数环面，泵送电荷$Q=eC$。体极化在一个周期内改变整数个电荷，却不要求瞬间产生可压缩金属态。Chern绝缘体与平凡绝缘体相邻时，两侧拓扑数差要求净手性边缘通道。可由Dirac质量畴壁直接求出：取$H=k\sigma_z+\ii\sigma_y\partial_y+m(y)\sigma_x$，若$m(y)$由负变正，零横向模满足一支$\psi(y)\propto\exp[-\int^y m(y')\dd y']$，另支指数增长不可归一。唯一可归一支沿边缘有$E=-k$。一支被舍掉的原因是空间归一化，不能将它当成另一个物理反向边缘态。
\originalsection{一般辛结构与相空间守恒}
课程的半经典框架还可脱离正则坐标来理解。设相空间变量为$\xi^I$，一阶作用量$L=a_I(\xi)\dot\xi^I-H(\xi)$，则变分给$b_{IJ}\dot\xi^J=\partial_IH$，$b_{IJ}=\partial_Ia_J-\partial_Ja_I$。$b$反对称且闭合$\partial_Ib_{JK}+\partial_Jb_{KI}+\partial_Kb_{IJ}=0$。若$b$可逆，$b^{IJ}$定义Poisson括号，其Jacobi恒等式来自这条闭合关系；任意指定一个反对称矩阵而不检查闭合性，未必是合法的哈密顿动力学。

相空间体积不是总为$\dd^{2d}\xi$，而是$\dd\Gamma=|\operatorname{Pf}b|\dd^{2d}\xi=\sqrt{\det b}\dd^{2d}\xi$，符号取定向后固定。辛流保持该体积：可由$\mathcal L_Xb=\dd(\iota_Xb)+\iota_X\dd b=\dd(\dd H)=0$推出其$d$次楔积也不变。因此占据函数$g$满足$\partial_tg+\dot\xi^I\partial_Ig=0$，等价的连续性式为$\partial_t(\sqrt{\det b}\,g)+\partial_I(\sqrt{\det b}\,g\dot\xi^I)=0$。Berry曲率导致的态密度修正正是这个测度的具体例子。

与固定温度环境耦合时，若只加扩散流$J^I=-\gamma^{IJ}\partial_Jg$，平衡趋向常数占据，对应无穷温；有限温需漂移配合。经典粒子可取$J^I=-\gamma^{IJ}g\partial_J(\log g+\beta H)$，使$g_0\propto\e^{-\beta H}$无耗散流。费米与玻色占据分别改为$J^I=-\gamma^{IJ}g(1\mp g)\partial_J[\log(g/(1\mp g))+\beta(H-\mu)]$。这里$\gamma$应有非负对称耗散部分，才能使相应自由能下降。这个开放体系模型保持粒子数却不要求体系自身能量守恒；封闭系统的碰撞项另须满足能量与动量守恒。
